\subsubsection{壮阳中药与西药的相互作用：同服之前必读}

"中药温和、西药猛烈，一起吃正好中和"是门诊反复听到的误解。中药同样有药理活性，与性功能常用西药同服时的相互作用值得单独列出：

\begin{itemize}
\item \textbf{淫羊藿（仙灵脾）}：其代表成分淫羊藿苷（icariin）在体外实验中表现出 PDE5 抑制样作用——这正是西地那非类药物的作用靶点。理论上与 PDE5 抑制剂同服存在叠加风险：血压下降、头痛、潮红、视觉异常；规律服用他达拉非者尤应主动告知医生；
\item \textbf{人参}：可能影响华法林等抗凝药的稳定性；红参的兴奋作用可能加重心悸，与降压药同用时血压波动需要监测；
\item \textbf{甘草}（多种复方中的常见成分）：长期大量使用可致假性醛固酮增多症——低血钾、水钠潴留、血压升高，与降压药、利尿剂的相互作用复杂，"甘草泡水养生"并非人人适宜；
\item \textbf{鹿茸及动物性"壮阳"制剂}：可能含微量激素样成分，与外源性睾酮等激素治疗叠加时，实际摄入剂量不可控。
\end{itemize}

\textbf{更要紧的是来源问题}：市场监管部门的抽检通报中，"纯中药壮阳药"被检出非法添加西地那非类似物的案例反复出现——这类产品剂量不明、标注不实，对正在服用硝酸酯类药物（硝酸甘油等）的冠心病患者可能致命（硝酸酯与 PDE5 抑制剂合用属绝对禁忌，可致严重低血压）。

因此三条原则：一、向医生如实告知你在用的所有中药、药酒与保健品——它们也是"药"；二、认准药品批准文号，拒绝"祖传秘方""三天见效"；三、正在服用 PDE5 抑制剂或硝酸酯类药物者，任何"壮阳"中药、药酒入嘴之前，先过一遍处方咨询。
